
\chapter*{Notational conventions}
\addcontentsline{toc}{chapter}{Notational conventions}
\markboth{Notational conventions}{Notational conventions}

When I refer to \textit{Japanese}, what I mean is \textit{Modern Standard Japanese}, the official version of the language spoken predominantly in the areas in and surrounding Tokyo. \textit{Premodern Japanese} is a broad term for any stage of the language prior to the Meiji period (1868–1912). \textit{Classical Japanese} (or \textit{bungo}), refers to the premodern (mostly) written stage of the language based on the literary style of the Heian-period (794–1185) and in use until the language reform carried out in the Meiji-period.

In the context of nominal modification, I will often prefer the term \textit{modifier} over \textit{adjective} with regard to Japanese. This is because adjectival status cannot always be unequivocally assigned to the modifiers under question. I will only use the term \textit{adjective} when I clearly refer to \textit{i}- and \textit{na}-adjectives respectively. With regard to other languages, I will in most cases stick to the terminology used by the respective authors.

I will represent syntactic trees for Japanese as head-final with the head to the right and the complement to the left. This means that in many cases no movement steps need to be assumed for Japanese that would be necessary if one follows to the bone the traditional cartographic system where everything is derived from a head-initial structure. Note, however, that this does not affect in any way the hierarchies and the structural precedences depicted here for the Japanese DP.

In the glosses, I will display copular elements, tense markers as well as the attributive \textsc{mod}-head, all of which are heads, connected to their preceding elements with a hyphen. This is done for ease of notation and readability. The tree structures will display them correctly as heads. Furthermore, I will present all monomoraic elements following Japanese modifiers as well as all copula forms in small caps throughout this book to ensure easy readibility.

All Japanese examples are transliterated following the Modern Hepburn Romanization System. I have chosen this system due to the fact that 1) pronunciation of Japanese material is depicted in a straightforward manner and 2) because the Modern Hepburn Romanization System is a unified and standardized system that does not allow room for ambiguity. This means that lengthened vowels, for instance [a:], will be displayed via a macron over the vowel ($<$\textit{ā}$>$). The only exception is [e:], which will be transliterated as $<$\textit{ei}$>$ sticking close to the original script. The semi-consonant /y/ will only be transliterated when it affects pronunciation, that is when it causes palatalization of the preceding consonant, not when it serves the function in the original Japanese script to display certain consonant-vowel combinations. For instance, the mora \textjpn{しゃ}, pronounced [\textctc a], will be transliterated as /sha/, which is very close to the original pronunciation, but never as $<$sya$>$ or $<$shya$>$. Furthermore, the object particle will be consistently transliterated as $<$\textit{wo}$>$. Transliterations of Persian follow the scientific DMG-system (DIN 31635), transliterations of languages with Cyrillic alphabets follow the scientific DIN-system as well.



